% Volume X: Normative Conclusions
% Part XVIII. Freedom Under Causal Opacity
% Chapter 106: The Right to Contest the Model
% Spine: proof-spines/ch106-spine.md

\chapter{The Right to Contest the Model}
\label{ch:ch106}

When an institution acts through a representation \(M(x)\), the represented person should ordinarily be able to know that the model exists, inspect the relevant inputs, challenge the associations it draws, introduce counterevidence, and seek revision. Model-governed power is not legitimate merely because it is technically sophisticated; it must remain contestable by the person it classifies.

The right to contest the model therefore reaches beyond a narrow complaint about outputs. It concerns the representational apparatus itself: what was treated as relevant, how signals were linked, which proxies were accepted, and by what procedure the model may be corrected. Where institutions govern through \(M(x)\), contestability is part of what makes representation politically tolerable.
